\begin{table}[H]
\centering
\caption{Treatment Effects on Lesson Completion Rate (Kenya)}
\label{tab:lesson_completion}
\begin{threeparttable}
\begin{tabular}[t]{l*{4}{>{\centering\arraybackslash}m{2.65cm}}}
\toprule
 & \multicolumn{4}{c}{Lesson completion rate} \\
\cmidrule(l{3pt}r{3pt}){2-5}
 & (1) & (2) & (3) & (4) \\
\midrule
Treat & 0.018 & 0.022 & 0.018 & 0.024 \\
 & [0.014] & [0.015] & [0.016] & [0.017] \\
\addlinespace
Teacher attendance &  & 0.085*** &  & 0.090*** \\
 &  & [0.029] &  & [0.034] \\
\addlinespace
Class size &  &  & -0.001 & -0.001 \\
 &  &  & [0.001] & [0.001] \\
\midrule
Control mean & 0.172 & 0.172 & 0.173 & 0.173 \\
Observations & 1128 & 1122 & 780 & 776 \\
\bottomrule
\end{tabular}
\begin{tablenotes}[para,flushleft]
\scriptsize
\setlength{\emergencystretch}{2em}
\item \textit{Notes:} Each column reports the coefficient from a regression of lesson completion rate on treatment status. The unit of observation is a school--subject--grade cell. Lesson completion is the fraction of scripted literacy-revision lessons completed, recorded in NewGlobe's classroom management system. Columns~2--4 add controls for teacher attendance rate and reading class size. Control means are computed in each column's estimation sample. Standard errors are clustered at the school level in brackets. ***, **, and * indicate significance at 1\%, 5\%, and 10\%.
\end{tablenotes}
\end{threeparttable}
\end{table}
